% Standalone solution dossier for ex:a5-gaussian-crossover.
\documentclass[../main.tex]{subfiles}
\begin{document}

% === SOLUTION HEADER (proof provenance) =====================================
% ledger-node: ex:a5-gaussian-crossover
% refines: ex:a5-gaussian-crossover
% bounded_by: none
% evidence_run: none
% checked_by: agent
% author: /root/a4_a5
% reviewer: /root/review_a4_a5
% dossier_editor: /root/review_a4_a5
% review: research/reviews/2026-08-21-a4-a5-independent-agent-audit.md
% date: 2026-08-21
% ============================================================================

\section*{Solution: exact Gaussian centred/non-centred crossover}

\begin{proposition}
Let $u,z\sim N(0,1)$ independently, $\theta=u+z$, and weight the law by the Gaussian likelihood
$\exp[-r(\theta-y)^2/2]$, $r\ge0$. In coordinates $(u,\theta)$ and $(u,z)$ the precision
matrices are
\[
 H_{\rm c}=\begin{pmatrix}2&-1\\-1&1+r\end{pmatrix},\qquad
 H_{\rm nc}=\begin{pmatrix}1+r&r\\r&1+r\end{pmatrix}.
\]
Both $C_P$ and the precision condition number favor non-centring for $0\le r<1$, tie at $r=1$,
and favor centring for $r>1$.
\end{proposition}

\begin{proof}
Expanding $u^2+(\theta-u)^2+r(\theta-y)^2$ and substituting $\theta=u+z$ gives the two matrices;
$y$ affects only the mean. Their determinants coincide and equal $1+2r$, while
\[
 \operatorname{tr}H_{\rm c}=3+r,\qquad
 \operatorname{tr}H_{\rm nc}=2+2r.
\]
For positive $2\times2$ matrices of fixed determinant, the smaller eigenvalue decreases and the
condition number increases with the trace. The trace difference is $1-r$, proving all three
regimes. Since a Gaussian with precision $H$ has $C_P=1/\lambda_{\min}(H)$, the same comparison
applies to $C_P$.
\end{proof}

\end{document}
